\chapter{Antecedentes}


\section{Fundamentos de CS}

\subsection{Descripción general del proceso}

Como se mencionó en la sección de introducción, el CS es una tecnología 
de fabricación que convencionalmente se utiliza para generar recubrimientos superficiales sobre componentes, mediante la proyección de partículas sólidas a alta velocidad sobre un sustrato. En la Figura \ref{CS} se presentan los elementos básicos que conforman un sistema estándar de CS.

La herramienta principal del proceso es una tobera convergente-divergente 
(o de Laval), la cual acelera un gas portador inerte (usualmente nitrógeno, helio, o mezclas de ambos) hasta velocidades supersónicas. Los polvos del material a depositar, generalmente metálicos, son suministrados por un alimentador y se mezclan con el flujo de gas en la tobera, siendo arrastrados y acelerados junto a este.

La velocidad alcanzada por el gas depende directamente de su temperatura 
tras el paso por el calefactor y de la presión de suministro. Existen 
configuraciones donde el polvo es precalentado de forma independiente antes de la inyección, con el fin de reducir su dureza y mejorar la probabilidad de adherencia; sin embargo, estas no constituyen la práctica convencional debido a la complejidad operativa que introducen.

Como se desarrolla en las secciones siguientes, la velocidad de impacto 
de las partículas es el factor determinante para la eficacia del proceso. 
Otros factores relevantes incluyen la morfología y el tamaño de las 
partículas, las propiedades mecánicas del material a depositar, y la 
naturaleza del sustrato.

El proceso se ejecuta desplazando la herramienta (o pistola) sobre la superficie a cubrir en sucesivas pasadas. La acumulación de material depende de parámetros operativos como el flujo másico de polvo, la velocidad de desplazamiento y el número de pasadas, entre otros.

Los depósitos obtenidos por CS suelen presentar porosidad residual, 
rugosidad superficial y tensiones residuales de compresión. Por esta razón, dependiendo de la aplicación, es frecuente aplicar tratamientos posteriores como recocido o mecanizado de acabado para mejorar las propiedades finales del recubrimiento.
 
\subsection{Aceleración del gas}

La tobera acelera el gas a velocidades supersónicas una vez que se supera el umbral establecido por la razón de presiones críticas entre el interior y exterior del sistema, generando el fénomeno de \textit{choking} o estrangulamiento, esto es:

\begin{equation}
	\frac{p_e}{p_0} < \left(\frac{2}{\gamma+1}\right)^{\frac{\gamma}{\gamma-1}}
	\label{eq:condicion_supersonica}
\end{equation}

donde \( p_e \) es la presión estática a la salida de la tobera (presión ambiental, por lo general), \( p_0 \) es la 
presión de estancamiento antes de la garganta (controlable), y \( \gamma \) es el coeficiente de calor específico del gas propulsor.

Es conveniente entonces caracterizar la velocidad de este flujo en función del número de Mach $M$:

\begin{equation}
	M = \frac{v}{a}
	\label{eq:mach}
\end{equation}

donde \( v \) es la velocidad local del flujo y \( a \) es la velocidad local del sonido en el gas, la cual se puede describir como:

\begin{equation}
	a = \sqrt{\gamma \frac{R}{M_{mol}} T}
	\label{eq:velocidad_sonido}
\end{equation}

donde, \( R \) es la constante universal de los gases ideales, 
\( M_{mol} \) es la masa molar del gas propulsor, y \( T \) es la temperatura estática local del flujo.

En la condición de estrangulamiento, el flujo siempre se encuentra en Mach 1 en la garganta, luego procede a expandirse, y acelerar, en la sección divergente. Asumiendo un flujo isoentrópico, la ecuación trascendental para encontrar el número de Mach a la salida es:

\begin{equation}
	\frac{A}{A^*} = \frac{1}{M} \left[ \frac{2}{\gamma+1} \left( 1 + \frac{\gamma-1}{2} M^2 \right) \right]^{\frac{\gamma+1}{2(\gamma-1)}}
\end{equation}

donde \( A \) es el área de la sección transversal local de la tobera, \( A^* \) es el área de la garganta en condición sónica (\( M = 1 \)), \( M \) es el número de Mach local del flujo, y \( \gamma \) es el coeficiente de calores específicos 
del gas propulsor. La cual suele resolverse de manera iterativa y da dos soluciones, donde la supersónica es la solución válida. 

Conocido el número de Mach, es posible encontrar la velocidad a la sálida mediante:

\begin{equation}
	v_e = M \sqrt{\frac{\gamma R T_0}{M_{mol}\left(1 + \frac{\gamma-1}{2}M^2\right)}}
	\label{eq:velocidad_salida}
\end{equation}

donde \( v_e \) es la velocidad del gas a la salida de la tobera, y \( T_0 \) es la temperatura de estancamiento antes de la tobera.

Así, este gas acelerado se mezcla en la tobera con los polvos de material a utilizar, otorgando energía cinética a las partículas a través de la transferencia de momentum lineal.

Nótese que en la ecuación~\ref{eq:velocidad_salida} la velocidad de salida del gas depende de la temperatura de estancamiento \( T_0 \), pero no de la presión de estancamiento \( p_0 \). Por lo tanto, una vez superada la condición de estrangulamiento, incrementar la presión de operación no produce un aumento en la velocidad del gas a la salida. Sin embargo, el aumento de presión sí incrementa la densidad del gas y, en consecuencia, el flujo másico a través de la tobera. Esto se traduce en una mayor tasa de transferencia de momentum lineal hacia las partículas, favoreciendo su aceleración. Así, la presión y la temperatura de operación actúan como parámetros independientes en el proceso: la temperatura determina la velocidad del gas, mientras que la presión regula la cantidad de momentum disponible para transferir a las partículas.

\subsection{Proyección de partículas}

El gas acelerado arrastra a las partículas en la tobera, aumentando su 
energía cinética mediante la transferencia de momentum lineal. Esta 
transferencia está determinada por varios factores relacionados con las 
condiciones de operación y las características del material particulado.

El punto de inyección de los polvos es uno de estos factores. Si se 
alimentan antes de la garganta, la distancia de interacción con el gas 
es mayor, lo que favorece que las partículas alcancen velocidades más 
altas. Sin embargo, esta configuración conlleva riesgos de obstrucción 
en la garganta por acumulación de material y de erosión acelerada de la 
tobera. Alimentar los polvos aguas abajo de la garganta evita estos 
problemas, pero a costa de reducir la longitud de interacción y, por 
tanto, la velocidad final de las partículas.

El tamaño de las partículas es otro factor determinante. En general, 
las partículas más pequeñas, al tener menor inercia, responden con mayor 
fidelidad a los cambios en el flujo de gas. Sin embargo, la relación entre 
tamaño y velocidad de impacto no es monótona: existe un tamaño óptimo para 
cada condición de operación. Las partículas muy finas, aunque se aceleran 
con mayor facilidad, son también más susceptibles de ser desviadas por 
fenómenos de turbulencia o por el \textit{bow shock} — el estancamiento 
del gas frente a la superficie del sustrato que genera una contrapresión 
localizada en la zona de impacto. Las partículas de mayor tamaño, por su 
inercia, son más capaces de superar esta barrera, aunque requieren mayor 
transferencia de momentum para alcanzar la velocidad crítica. Por tanto, 
la selección del rango granulométrico implica un compromiso entre 
acelerabilidad y capacidad de penetrar la capa de gas estancado frente 
al sustrato.

Comercialmente, los polvos para CS se describen y adquieren según su 
distribución de tamaño de partícula, caracterizada mediante percentiles 
del diámetro equivalente: $D_{10}$, $D_{50}$ y $D_{90}$, que corresponden 
a los diámetros bajo los cuales se encuentra el 10\%, 50\% y 90\% del 
volumen de polvo, respectivamente. En CS interesa principalmente el rango 
$D_{10}$--$D_{90}$, pues define la fracción de partículas con mayor 
probabilidad de encontrarse dentro de la ventana entre la velocidad crítica 
y la velocidad de erosión del depósito.

El tamaño de partícula también influye en su calentamiento por convección 
durante el trayecto en la tobera: las partículas más pequeñas equilibran 
su temperatura con la del gas más rápidamente. A diferencia de otros 
procesos de proyección térmica, en CS las partículas no alcanzan su punto 
de fusión, por lo que la temperatura de la partícula no es el factor 
dominante del proceso. No obstante, no es un factor irrelevante: una 
partícula más caliente presenta menor dureza y mayor ductilidad, lo que 
reduce su velocidad crítica de adhesión. En este sentido, la temperatura 
del gas de proceso actúa sobre la eficacia del CS por dos vías: 
incrementando la velocidad del gas —y con ello la velocidad de impacto 
de las partículas— y aumentando la temperatura de las partículas, 
favoreciendo su deformabilidad en el impacto.

Finalmente, la morfología de las partículas también es relevante. 
Las partículas esféricas, producidas típicamente por atomización en gas, 
presentan un comportamiento aerodinámico más predecible y una mejor 
fluidez en el alimentador, siendo la morfología preferida en CS 
convencional. Existe interés en partículas de forma irregular en 
contextos específicos, aunque los efectos sobre la calidad del depósito 
son dependientes del material y las condiciones de operación [REFERENCIAS]


\subsection{Impacto}

\subsection{Velocidad crítica}

\subsection{Mecanismos de adherencia}

Adiabatic shear inestability asdasd

Anclaje mecánico

Difusión metálica

Enlace intermetálico

Fuerzas de van der Wals

De lo anterior, es claro que entonces debe existir una velocidad mínima a la cual ocurra suficiente deformación plástica, pues de lo contrario, la partícula podría deformarse (de manera plástica y/o elástica) pero no lo suficiente para adherirse. A esta velocidad se le llama \textbf{velocidad crítica}. 

velocidad critica

velocidad maxima

tamaño de particula

El cambio de energpia interna en las partículas por contacto con el gas de proceso no suele ser relevante en el proceso de adhesión, pues las temperatura del gas suele ser basntante lejana a las temperaturas de cambio de fase o recristalización de los materiales típcios de coldspray, considerando además que el tiempo durante el cual la partícula se mezcla con el gas es reducido. La temperatura podría afecta la ductilidad de las par´ticuñas y facilitar la deformación 

\subsection{Perfil de deposición}

\section{Fabricación por CS}

\subsection{Sistema de CS}
\subsection{Parámetros relevantes}
\subsection{Características de productos por CS}
\subsection{Aplicaciones típicas}

\section{Manufactura aditiva}
\subsection{Slicing y trayectorias}
\subsection{Estrategias de deposición}
\subsection{Control multieje}
\subsection{FA usando CS}
\subsection{FA usando titanio}


